% ============================================================
%  扉页
% ============================================================
\begin{titlepage}
\thispagestyle{empty}
\centering
\vspace*{0.8cm}

{\small\color{shuyiGray}数译 · 威腾演讲译注　第一种}

\vspace{2.0cm}
{\color{shuyiBlue}\rule{0.74\textwidth}{0.9pt}}

\vspace{1.0cm}
{\Huge\bfseries 对偶、Langlands 对应\\[0.42em]与 Michael Atiyah 的影响\par}

\vspace{1.0cm}
{\Large Edward Witten\par}
\vspace{0.5cm}
{\normalsize\color{shuyiGray}2021 年 9 月 23 日，Isaac Newton Institute 演讲\par}

\vspace{0.7cm}
{\color{shuyiBlue}\rule{0.74\textwidth}{0.9pt}}

\vfill

\begin{minipage}{0.76\textwidth}
\centering\small\color{shuyiGray}
从 1977 年春天 Atiyah 的一句提醒，到 2024 年几何 Langlands 的证明，
一条横跨量子场论与数论的暗线，在四十七年后合拢。
\end{minipage}

\vfill

{\large 中文译注本\par}
\vspace{0.45cm}
{\large\bfseries 超对称梦想家　译注\par}

\vspace{2.2cm}
{\small 2026 年 11 月\par}

\vspace{0.6cm}
{\scriptsize\color{shuyiGray}
底本：\emph{Duality, the Langlands Correspondence, and the Influence of Michael Atiyah}\\
54 张 beamer 幻灯片（16:9），Edward Witten，2021 年 9 月 23 日\par}

\end{titlepage}
